\documentclass[10pt]{article}
\usepackage[margin=20mm]{geometry}
\usepackage{amsmath,amssymb,mathtools,fontspec,unicode-math,xurl,hyperref}
\IfFontExistsTF{DejaVu Serif}{\setmainfont{DejaVu Serif}}{\setmainfont{Latin Modern Roman}}
\setmathfont{latinmodern-math.otf}
\setlength{\parindent}{0pt}\setlength{\parskip}{6pt}\sloppy
\title{Initial Stability for Ricci Flow: A Claim-Focused Review}
\author{AI-assisted technical review}
\date{}
\begin{document}
\maketitle
Abstract: This review checks the local stability mechanism and its application to uniformly Ricci-pinched three-manifolds in the acquired author manuscript of Deruelle, Schulze, and Simon. The checked scaling, localization, excess, and rigidity arguments support the qualitative conclusions within the stated hypotheses. No substantive false theorem is established. One precise quantifier in the higher-derivative conclusion remains unresolved: the printed common time interval with no order-dependent prefactor is stronger than the fixed-order interpolation argument explicitly supplies. The assessment at this point is Needs Evidence, not a proved counterexample. Scientific adjudication remains pending.\par
\section*{Paper identity and scope}
The paper is A. Deruelle, F. Schulze, and M. Simon, Initial stability estimates for Ricci flow and three dimensional Ricci-pinched manifolds [DSS]. The exact reviewed source is arXiv:2203.15313v2, submitted 17 March 2025, with 50 PDF pages. Its associated journal DOI is https://doi.org/10.1215/00127094-2025-0018. The publisher PDF was not acquired; no identity of its text with this manuscript is asserted. Page references below are one-based PDF pages, which agree with the manuscript’s printed page numbers.\par
The review covers the statements and dependencies of Theorems 1.2 and 1.3, the curvature and coordinate adjustments in Sections 2–4, the polynomial and exponential estimates in Sections 5–6, the rigidity/application chain in Sections 7–8, and the scaling and splitting mechanisms in Appendices A–B. It is a claim-focused initial review, not an independent reconstruction of every imported compactness, pseudolocality, Ricci-limit, or Alexandrov theorem. The coverage ledger distinguishes direct checks, bounded native AI assistance, and imported results not reproved. No computational experiment is needed for the mathematical statements under examination.\par
\section*{Hypotheses and dependencies}
The stability setting consists of smooth local Ricci flows for positive time, with \(|\operatorname{Rm}(g(t))|\leq c_0/t\) and \(\operatorname{Ric}(g(t))\geq-g(t)\), a common limiting distance up to isometry, Euclidean tangent cones, and sufficiently close-to-isometric distance coordinates. Completeness is not imposed on the local flows; the relative compactness assumptions used in the local distance and curvature arguments must be retained. The conclusion compares flows after harmonic-map gauges and on a smaller ball and time interval. An arbitrary comparison of ungauged metric components is not a test of this existential conclusion.\par
For Theorem 1.3, uniformly Ricci pinched explicitly includes \(\operatorname{Ric}\geq0\) together with \(\operatorname{Ric}\geq c\,\mathrm{R}g\) for some \(c>0\). A negatively curved Einstein metric does not satisfy this definition. Completeness, dimension three, and bounded initial curvature belong to the stated theorem. The higher-dimensional questions are not proved extensions.\par
\section*{Checks supporting the argument}
\subsection*{Scaling and the initial gauge}
Under \(\widehat g(t)=\lambda^2g(t/\lambda^2)\), \[t|\operatorname{Rm}(\widehat g(t))|
=(t/\lambda^2)|\operatorname{Rm}(g(t/\lambda^2))|.\] The lower Ricci bound improves by \(\lambda^{-2}\) in the rescaled metric. These are the scales required in Appendix A. The fixed-domain volume inequality used there also has the correct direction: if \(\mathrm{R}\geq-\eta\), then for a fixed measurable \(\Omega\) and \(r<s\), \[\operatorname{Vol}_{g(s)}(\Omega)
 \leq e^{\eta(s-r)}\operatorname{Vol}_{g(r)}(\Omega).\] A near-Euclidean lower bound at the later time therefore supplies the needed earlier-time lower bound.\par
The adjustment and harmonic-map steps precede comparison of metric tensors. Their scale-dependent derivative estimates and shrinking interior domains are consistent with the energy argument. A distance-space limit alone was not treated as tensor convergence. Literal index and constant slips were checked against subsequent use; no main-conclusion defect follows from those slips alone.\par
\subsection*{Polynomial and exponential decay}
Write \(h=\widetilde g_2-\widetilde g_1\) and \(u=t^{-C}|h|_\delta^2\). The Section 5 estimates supply decay faster than any fixed power after the required domain restriction. Thus the initial term vanishes when the localized divergence-form inequality for \(u\) is integrated from a positive time and that time tends to zero. The cutoff source is bounded and supported away from the observation ball.\par
For \(|x|<R_0/4\) and \(R_0/2\leq |y|<R_0\), the safe separation estimate is \(|x-y|\geq R_0/4\). The relevant Gaussian integral is therefore bounded by \[C\int_0^t s^{-n/2}\exp\!\left(-\frac{aR_0^2}{16s}\right)\,ds .\] Substitution \(\tau=R_0^2/s\) gives a polynomial times an exponential tail. Fixed polynomial prefactors can be absorbed by reducing the positive exponential constant and time interval. The printed factor \(3/16\) on page 29 is not the geometric lower bound for the stated balls; replacing it by \(1/16\) preserves Theorem 6.2, whose exponent is unspecified. This establishes no failure of its qualitative rate.\par
\subsection*{Splitting and the application}
The product-space excess estimate in Appendix B admits a direct check. Put \(a=\delta/\varepsilon\), \(b=1/\delta\), and suppose \(R\geq20(a+b)\). The longitudinal distances satisfy \(s_\pm\geq19R/20\) and the transverse distances satisfy \(d_\pm\leq a+b\). Dropping the nonnegative transverse contribution from the subtracted long distance gives \[0\leq\mathcal E_R
 \leq\sum_\pm\bigl(\sqrt{s_\pm^2+d_\pm^2}-s_\pm\bigr)
 \leq\frac{40}{19R}(a^2+b^2).\] This gives the required vanishing excess with the selected \(R_i\). The opposite limiting Busemann functions are harmonic. Bochner’s formula and the maximum principle give unit parallel gradient; its contraction with Ricci vanishes. Since \(\partial_tg=-2\operatorname{Ric}\), the gradient vector field is independent of time, allowing a common product identification.\par
The compactness application depends on noncollapse, the three-dimensional Ricci-limit/cone structure, the comparison expander, and the almost-pinched rigidity result. Scaling from small time on an annulus to decay at infinity has been checked. The comparison transfers an exponentially small pinching error; it does not assert exact equality of the flows. The external asymptotic criterion in Deruelle’s Theorem 1.3 [Deruelle] was checked for applicability: \(f\asymp r^2\) and \(|\operatorname{Ric}|=O(e^{-f^\beta})\), \(\beta>0\), imply \(r^2|\operatorname{Ric}|\to0\), as that theorem requires. Its complete proof, and the other imported existence and structure results, were not independently reconstructed.\par
\section*{Unresolved quantifier}
Equation reference: eq:printed. Problem 1 (A common interval for every derivative order).  Theorem 1.2, page 4, equation (1.7), states that for every \(k\geq0\) there is \(C_k>0\) such that \[
 |D^kh(t)|_\delta\leq e^{-C_k/t}
 \quad\text{for all }t\in(0,T_0],\] with the same \(T_0\) introduced earlier. The proof on page 30 invokes interpolation and the previously obtained \(C^0\) estimate.\par
Equation reference: eq:available. The fixed-order inference directly gives a bound of the form \[
 |D^kh(t)|_\delta\leq A_k t^{-m_k}e^{-a_k/t},
 \qquad A_k,a_k>0,\] on a suitable smaller spatial region. For each fixed \(k\), a reduced time interval absorbs \(A_kt^{-m_k}\). This does not itself show that those intervals contain a common positive time interval across every order.\par
A diagnostic of this logical distinction is \[f(x,t)=\frac{e^{-1/t}}{1-x},\qquad |x|<\tfrac12,\quad t>0.\] Each fixed derivative decays exponentially as \(t\downarrow0\), whereas \(D_x^kf(0,t)=k!e^{-1/t}\) is unbounded in \(k\) at each fixed positive time. Thus fixed-order exponential estimates do not generally imply [eq:printed] (\#eq:printed), whose right-hand side is less than one. This scalar function is not a Ricci-flow counterexample and does not account for the theorem’s freedom of gauge.\par
The missing evidence is an argument controlling all derivative orders on the common interval, or clarification that the intended interval may depend on \(k\). An order-dependent prefactor as in [eq:available] (\#eq:available), or separate \(T_k>0\), is consistent with the checked interpolation step. This observation does not disprove the geometric theorem, does not establish a Major or Minor error, and does not undermine the finite-order estimates needed for the curvature comparison in Section 8.\par
\section*{Assessment}
No substantive counterexample to local \(C^0\) exponential stability or the three-dimensional compact-or-flat conclusion was established. The exact all-orders/common-time phrasing remains Needs Evidence. Displayed sign and constant discrepancies were separated from demonstrated effects on the conclusions; no severity quota was used.\par
Dimension | State\par
Inspection | Needs Evidence for the derivative quantifier; no confirmed substantive defect in the other declared checks.\par
Repair verification | Fixed-order prefactor absorption checked; the stronger common-interval claim remains unproved here.\par
Author adoption | Not requested; no adoption or author correction asserted.\par
Scientific decision | Pending; no human review recorded.\par
Remaining work | Resolve the common-interval quantifier; compare the publisher text if that exact version becomes available.\par
\section*{Evidence and limitations}
Analytic checks are distinct from numerical tests, author-code execution, and empirical reproduction; none of the latter is claimed. The source PDF, acquired TeX, coverage, native-helper notes, source-page images, exact report build, retained dependencies, and page-view receipts are bound by the verification records. The source PDF SHA-256 is `89e8f91db4ffe05e82b48384c73db7c24ed9ab59bcd74b330201cd92ba64f8c6`. This is an AI initial review with bounded native AI assistance, not independent acceptance or a human attestation. No equivalence to unacquired publisher material, and no universal absence of supplements or later corrections, is asserted.\par
2 A. Deruelle, F. Schulze, and M. Simon, Initial stability estimates for Ricci flow and three dimensional Ricci-pinched manifolds, arXiv:2203.15313v2, 17 March 2025. https://arxiv.org/abs/2203.15313v2. A. Deruelle, Asymptotic estimates and compactness of expanding gradient Ricci solitons, author-hosted 39-page PDF snapshot retrieved 26 October 2026; Theorem 1.3 on page 3 and Proposition 3.5 discussion on page 24 consulted by the bounded native helper. https://perso.imj-prg.fr/alix-deruelle/wp-content/uploads/deruelle-pub/Comp-Asy-EGS.pdf.\par
\end{document}
